\section{讨论与结论}
\label{sec:conclusion}
本文围绕生成式三维资产缺少部件身份的问题，将视觉语言模型的作用落实为可检查的空间提示。Qwen3-VL 负责可见部件规划与候选反馈，SAM2 负责掩码预测，点引导清理和小框优先仲裁负责形成互斥二维引导，冻结 SegViGen 负责后续三维分割。英文名称通过部件记录保留，避免将“已有名称”“已有二维区域”和“已有三维语义绑定”混为一谈。该设计不需要为当前任务重新训练这些基础模型，但其计算成本仍包括语言规划、反馈与三维推理。

方法的局限首先来自单视角可见性：被遮挡部件无法仅靠源图验证，二维投影分数也无法覆盖完整表面。其次，视觉语言规划可能遗漏小件、重复命名或提出错误负点；小框优先融合只解决冲突分配，不能保证每个候选在语义上成立。残余像素填充能够补齐前景，却不会自动发现新的部件。最后，当前保存的语义记录与二维标签明确对应，但三维表面到部件 ID 的持久映射仍需实现与单独验证。

现有材料支持形成可复现的方法描述和实验方案，尚不足以断言本文方法优于自动 SAM2 或其他三维分割方法。后续将按冻结测试集补齐主比较、消融和效率结果，并据此收敛最终结论；同时完善三维 ID 导出与交互展示验证，使名称不仅存在于记录中，也能被可靠地查询到具体三维区域。
